\documentclass[12pt]{article}
\usepackage{amsmath,amssymb}

\begin{document}
\section*{{2025 AIME II Problems/Problem 9}}
Source: AoPS Wiki

\subsection*{Solution}
For $\sin(7\pi\cdot\sin(5x))=0$ to happen, whatever is inside the function must be of form $k\pi$ . We then equate to have
\begin{align*}
7\pi\cdot\sin(5x)=k\pi \\
\sin(5x)=\frac{k}{7} \\
\end{align*}
We know that $-1\le \sin{5x} \le 1$ , so clearly $k$ takes all values $-7\le k \le 7$ . Since the graph of $\sin{5x}$ has 5 periods between $0$ and $360$ , each of the values $k=-6,-5,-4...-1,1,2...6$ give $10$ solutions each. $k=-7,7$ give $5$ solutions each and $k=0$ gives $9$ solutions (to verify this sketch a graph). Thus, $n=139$ . We know that the function is tangent to the x-axis if it retains the same sign on both sides of the function. This is not true for points at $k=-6,-5,-4...4,5,6$ because one side will be positive and one will be negative. However this will happen if $k=-7,7$ because the sine function "bounces back" and goes over the same values again (to clarify, because $y=\sin{5x}$ has the same values on both sides of $\sin{5x}=\frac{k}{7}$ for $k=-7,7$ , as it's a maximum $\sin{5x}=\frac{7}{7}=1$ or minimum $\sin{5x}=\frac{-7}{7}=-1$ , the corresponding values of $y=\sin(7\pi\cdot\sin(5x))$ also have a maximum/minimum and repeat around those $x$ -values), and $t=10$ of these values exist.
Thus, $n+t=\boxed{149}$ . ~ lisztepos ~ Clarification by Christian

\end{document}
